\section{Mở rộng sang video trên PVSG}
\subsection{Từ một cạnh trong ảnh đến một quan hệ có thời hạn}
Ở thiết lập PredCls trên ảnh tĩnh, mô hình có thể dự đoán một quan hệ cho cặp
đối tượng từ một lần quan sát. Trong video, hai đối tượng có thể xuất hiện cùng
nhau lâu hơn hành động liên kết chúng. Chẳng hạn, người và chiếc cốc có thể
cùng ở trong khung hình cả trước khi người cầm cốc và sau khi đặt cốc xuống.
Đầu ra vì thế gồm ID chủ thể $s$, predicate $p$, ID đối tượng $o$ và tập frame
$A_{spo}$ trên đó quan hệ có hiệu lực. Một bộ ba có thể có nhiều đoạn rời nhau.
Hai cặp có thứ tự $(s,o)$ và $(o,s)$ được chấm riêng.

PVSG tách việc tìm và theo dõi vùng đối tượng khỏi việc tìm quan hệ trên
các track~\cite{yang2023pvsg}. Thực nghiệm này cấp mask, nhãn lớp và ID chuẩn
của đối tượng ở từng frame. ID nối các quan sát của cùng một đối tượng qua
thời gian. Mô hình phải tự chọn cặp, predicate và khoảng thời gian. Thiết lập
này cho phép quy khác biệt giữa các cấu hình cho phần dự đoán quan hệ khi đầu
vào đối tượng như nhau. Phạm vi khác với kiểm tra bảo toàn nội dung sau khi
chỉnh sửa video trên DAVIS hoặc YouTube-VOS được gợi ý trong đề.

\subsection{Dữ liệu, căn chỉnh và biểu diễn}
Bản PVSG phát hành gồm 400 video tại 5 frame/giây với 149.488 frame chú thích,
126 lớp đối tượng và 57 predicate. Ba nguồn là VidOR (289 video), Ego4D (56)
và EPIC-Kitchens (55). Nhóm giữ danh sách 338 video train và 62 video test
của bản phát hành. Tập train được chia thành 287 video fit và 51 video
validation. Việc chia theo video giữ các cửa sổ chồng lấn của một video trong
cùng một tập. Có bốn nhóm ID nguồn Ego4D xuất hiện ở cả hai phía của split phát
hành. Vì vậy độ độc lập theo cảnh quay gốc còn là một giới hạn.

Mask xác định trục thời gian và vùng nhìn thấy của mỗi track. Kiểm tra media phát
hiện bốn frame RGB thiếu; tại những vị trí này vẫn giữ thông tin hình học và
cờ chỉ báo thiếu đặc trưng thị giác. Với nhãn quan hệ, nhóm lấy giao thời gian
chú thích với các frame hai track cùng nhìn thấy. Các nhãn ngoài từ vựng,
tham chiếu track sai và đoạn rỗng sau phép giao được ghi vào nhật ký xử lý.
Cách làm này giữ nguyên danh sách video và cho mọi mô hình cùng một tập đích.
CLIP ViT-B/32~\cite{radford2021clip} đóng băng trích vector ảnh 512 chiều; cache dùng chung cho cả
400 video. Nhờ vậy có thể thay relation head mà vẫn dùng đúng đặc trưng
của cùng một vùng trên cùng một frame.

\subsection{Đồ thị và các mô hình đối chứng}
Một đỉnh biểu diễn một đối tượng nhìn thấy tại một frame. Cạnh không gian
cho phép trao đổi thông tin giữa đối tượng cùng frame. Cạnh thời gian nối hai
quan sát liền kề của cùng track. Tập cặp cần dự đoán gồm mọi cặp có hướng,
khác ID, có ít nhất một frame đồng hiện. Tập cặp này vẫn giữ nguyên khi bỏ
cạnh truyền thông điệp trong mô hình. Nhờ đó ablation đo ảnh hưởng của ngữ
cảnh mà không làm bài toán dễ đi bằng cách loại ứng viên.

MLP mã hóa đối tượng và chấm từng cặp ở từng frame. Static GAT cho phép các
đỉnh lấy ngữ cảnh trong một frame trước khi chấm cặp. ST-GAT thêm ngữ cảnh
trên track; hai lớp thời gian cho trường tiếp nhận tối đa năm frame.
Ba biến thể lần lượt bỏ cạnh thời gian, bỏ cạnh không gian và bỏ cả hai loại
cạnh giữa các đỉnh. Cả ba giữ self-loop và cùng kiến trúc ST-GAT. Do đó Static
GAT và ST-GAT bỏ cạnh thời gian vẫn là hai cấu hình riêng.

Mô hình Temporal Conv xử lý chuỗi đặc trưng của cùng một cặp đối tượng. Khi một
đối tượng khuất, chuỗi được tách tại chỗ mất đồng hiện. Conv kernel 5, ba lớp
chuyển thể thiết kế trong mục 5.2 của bài PVSG. Nhóm dùng projection cặp và
đầu ra đa nhãn chung thay cho các đầu ra riêng của OpenPVSG. Conv kernel 3,
hai lớp là đối chứng bổ sung. Trường tiếp nhận tối đa tương ứng là 13 và 5 frame.
Trong bài CVPR, segmentation và tracking cung cấp track dự đoán, nên điểm
end-to-end còn phản ánh cả lỗi mask, lớp và liên kết ID. Bảng của nhóm chỉ so
sánh trực tiếp tám cấu hình dùng cùng GT track.

\subsection{Từ điểm từng frame đến khoảng dự đoán}
Với mỗi cặp và frame hợp lệ, mô hình xuất 57 logits; masked binary
cross-entropy cho phép nhiều quan hệ đồng thời. Khi đánh giá, cửa sổ dài 64
frame, bước trượt 32. Với frame thuộc nhiều cửa sổ, lấy trung bình logits
trước sigmoid rồi chọn các frame có xác suất ít nhất 0,5. Mỗi đoạn frame liên
tiếp tạo một interval với hai đầu mút tính cả. Điểm xếp hạng của bộ ba là xác
suất trung bình trên các frame đã chọn. Cách này trả về nhiều khoảng cho cùng
bộ ba và xếp hạng sau khi đã ghép toàn bộ video.

Gọi $A$ là tập frame dự đoán và $B$ là tập frame chuẩn của cùng bộ ba.
Khi đó $\mathrm{tIoU}=|A\cap B|/|A\cup B|$. Ví dụ, nếu nhãn gồm các
frame 10--14 và dự đoán gồm các frame 12--16, phần giao có ba frame và
phần hợp có bảy; bộ ba đúng vẫn bị tính trượt ở ngưỡng 0,5. Khi $B$ có hai
đoạn rời, phép tính trên tập frame giữ được khoảng trống giữa hai đoạn.

Mỗi video lấy tối đa 50 bộ ba có điểm cao nhất. Một lượt trúng cần khớp
đúng thứ tự hai track, predicate và đạt tIoU ít nhất 0,5. R@50 lấy tổng
số trúng chia cho 931 bộ ba chuẩn trên 62 video. mR@50 tính recall riêng
từng lớp rồi lấy trung bình trên 54 lớp có mặt trong test. Cách gộp này khác
với trung bình trên từng ảnh trong protocol VG150.

\subsection{Huấn luyện và chọn checkpoint}
Mỗi cấu hình chạy ba seed 17, 42 và 2026 với batch 32, tối đa 150 epoch
và patience 15. Tiêu chí chọn là mR@50 validation, hòa thì xét R@50 rồi
lấy epoch sớm hơn. Sau khi xem validation và trước khi chấm test, nhóm
áp dụng cùng một điều kiện kiểm tra hội tụ muộn cho toàn bộ 24 cặp mô hình/seed.
Nếu lượt gốc dừng ở epoch 16 với cả hai recall validation bằng 0, huấn
luyện lại từ đầu với cùng seed và chỉ cho phép dừng từ epoch 50.
Bảng~\ref{tab:pvsg} gồm 17 lượt gốc và 7 lượt huấn luyện lại; mỗi seed
đóng góp đúng một giá trị. Cách chọn bảo đảm cùng tiêu chí giữa các seed
trong khi cho phép chi phí huấn luyện thực tế khác nhau. Bảy so sánh từng
cặp nằm ở Phụ lục~\ref{app:late}.

\begin{table}[!ht]\centering
\caption{PVSG test trên track chuẩn. Bảng chính có tính hội tụ muộn;
trung bình $\pm$ độ lệch chuẩn mẫu của ba seed, đơn vị \%.}\label{tab:pvsg}
\small\setlength{\tabcolsep}{3pt}
\begin{tabular}{@{}lrr@{}}\toprule
Mô hình & R@50 & mR@50\\\midrule
MLP & $3,80\pm0,43$ & $1,62\pm0,12$\\
Static GAT & $3,80\pm0,65$ & $1,68\pm0,17$\\
ST-GAT & $4,98\pm0,27$ & $\mathbf{2,38\pm0,54}$\\
ST-GAT bỏ cạnh thời gian & $3,58\pm1,35$ & $1,28\pm0,75$\\
ST-GAT bỏ cạnh không gian & $3,47\pm1,20$ & $1,07\pm0,70$\\
ST-GAT chỉ giữ self-loop & $3,62\pm0,16$ & $1,51\pm0,04$\\
Conv kernel 3, hai lớp & $6,77\pm0,65$ & $1,29\pm0,26$\\
Conv kernel 5, ba lớp & $\mathbf{8,63\pm0,61}$ & $1,25\pm0,25$\\
\bottomrule\end{tabular}
\end{table}

\subsection{Recall tổng, recall theo lớp và sai số thời gian}
Conv kernel 5 thu hồi nhiều bộ ba nhất; recall theo nguồn Ego4D,
EPIC-Kitchens và VidOR lần lượt là 11,72\%, 10,67\% và 7,18\%. Khi xếp 57
predicate theo tần suất trên tập fit và chia đều thành ba nhóm, mR@50 của
Conv kernel 5 lần lượt là 3,56\%, 0,00\% và 0,00\% từ phổ biến đến hiếm.
Một mô hình tập trung vào lớp phổ biến vì thế có thể tăng R@50 dù nhiều lớp
chưa có lượt trúng. ST-GAT đầy đủ đạt mR@50 cao nhất và có mean cao hơn cả
ba ablation cạnh. Với ba seed và độ phân tán quan sát được, đây là nhận xét
về các lượt đã đo. Hai Conv khác cả độ sâu lẫn số tham số; mức tăng R@50
của Conv kernel 5 chưa tách riêng được tác dụng của cửa sổ ngữ cảnh.

Ở Static GAT seed 17, bộ ba $(8,20,4)$ trong video \path{0004_11566980553}
được xếp trong top 50 nhưng đạt tIoU 0,289. Ở video
\path{0039_6951351121}, bộ ba $(14,43,3)$ có tIoU đúng bằng 0,50 và được
tính trúng. Hình~\ref{fig:timeline} đối chiếu một số tập frame chuẩn và dự
đoán. Cả hai ví dụ cho thấy phân tích chỉ ở mức tên quan hệ dễ bỏ qua sai số
định vị thời gian.
